\documentclass[11pt]{article}
\usepackage[margin=1in]{geometry}
\usepackage{amsmath,amssymb}
\usepackage{hyperref}
\usepackage{booktabs}
\usepackage{xcolor}
\usepackage{listings}

\lstset{
  basicstyle=\ttfamily\footnotesize,
  breaklines=true,
  columns=fullflexible,
  showstringspaces=false,
  xleftmargin=1em,
  frame=single,
  framesep=3pt,
}

\title{Update 5: Upstream UCX Pull Request Review, External Validation,
  and Middleware Hardening\\[4pt]
  \large Post-Submission Follow-Up to Update 4}
\author{Nathan Gopee\\MS Computer Science Candidate, SUNY New Paltz\\\texttt{gopeen1@newpaltz.edu}\\[4pt]Advisors: Dr.\ Wafi Danesh \& Dr.\ Ashley Suchy}
\date{April 12, 2026}

\begin{document}
\maketitle
\tableofcontents
\newpage

\section{Recap}

Update 4 introduced libscuffedrdma's dual-QP architecture, the PMP
bang-bang controller, the WFA classifier, and the landscape of upstream
UCX issues that motivated six pull requests (\#11304--\#11309). This
update collects what happened after those PRs were submitted:

\begin{itemize}
  \item The actual code diffs landed on each branch and their current
        commit hashes.
  \item External validation that arrived in the same month, most
        importantly the \texttt{autoscriptlabs/libmesh-rdma} project and
        its accompanying r/HPC thread.
  \item Two contrasting designs from the production LLM systems space
        (Apple MLX PR \#3094 and Pure Storage TurboQuant + FlashBlade).
  \item A security audit of all six PRs and of the libscuffedrdma
        middleware, triggered by a memory-safety bug the author found
        in libmesh-rdma's \texttt{mesh\_mpi\_pt2pt.c}.
  \item A set of hardening commits applied to libscuffedrdma's Python
        layer, adapting three modules from libmesh-rdma while
        remediating four high-severity middleware findings.
\end{itemize}

Update 4 stays as-is; everything below is additive.

\section{Pull Request Review with Real Diffs}

Each PR is summarised below with the exact fix landed on the branch.
All hashes come from the \texttt{ndg8743/ucx} fork at the time of
writing.

\subsection{PR \#11304 --- Calculated rndv threshold for \texttt{ucp\_tag\_send\_nbr()}}
\label{sec:pr_11304}

Fixes \href{https://github.com/openucx/ucx/issues/4430}{\#4430}. Head
commit \texttt{694b20db8}. Default value flipped to \texttt{auto} so
both send paths use the transport-aware threshold:

\begin{lstlisting}[language=C]
-  {"RNDV_SEND_NBR_THRESH", "256k",
+  {"RNDV_SEND_NBR_THRESH", "auto",
   "Threshold for switching from eager to rendezvous protocol in "
   "ucp_tag_send_nbr().\n"
+  "\"auto\" means to use the same dynamically calculated threshold "
+  "as for regular sends.\n"
\end{lstlisting}

With the ``auto'' sentinel in place, the rendezvous threshold derivation
falls through to \texttt{rndv\_thresh} (the calculated value) rather than
the hardcoded 256 KB:

\begin{lstlisting}[language=C]
- rndv_local_thresh = context->config.ext.rndv_send_nbr_thresh;
+ if (context->config.ext.rndv_send_nbr_thresh == UCS_MEMUNITS_AUTO) {
+     rndv_local_thresh = rndv_thresh;
+ } else {
+     rndv_local_thresh = context->config.ext.rndv_send_nbr_thresh;
+ }
\end{lstlisting}

This is a pure behavior-flip: no new code paths, no new allocations,
and no change to any on-wire protocol. Explicit values
(\texttt{UCX\_RNDV\_SEND\_NBR\_THRESH=256k}) still work for users who
want the old behavior.

\subsection{PR \#11305 --- Adaptive TX CQ moderation}
\label{sec:pr_11305}

Fixes \href{https://github.com/openucx/ucx/issues/1307}{\#1307}. Head
commit \texttt{e1e419982}. Three lines in one inline function:

\begin{lstlisting}[language=C]
-    return (txqp->unsignaled >= iface->config.tx_moderation) ? flag : 0;
+    return (txqp->unsignaled >= iface->config.tx_moderation) ||
+           (iface->tx.cq_available <= (signed)iface->config.tx_moderation)
+           ? flag : 0;
\end{lstlisting}

Without the second disjunct, fan-out workloads where each RC endpoint
sends fewer than \texttt{tx\_moderation} messages never produce a
signaled send. The TX CQ fills up with unsignaled completions, credits
are never reclaimed, and the buffer pool deadlocks. Adding the
iface-level \texttt{cq\_available} check forces the next send on any
endpoint to be signaled once CQ credits drop to the moderation
threshold.

\subsection{PR \#11306 --- Eager inline sends for host memory when CUDA MDs are present}
\label{sec:pr_11306}

Fixes \href{https://github.com/openucx/ucx/issues/4275}{\#4275}. Head
commit \texttt{39939fb4e}. Adds a memory-type cache lookup to
\texttt{ucp\_proto\_is\_inline}:

\begin{lstlisting}[language=C]
+  /* Look up the buffer in the memory type cache. If not found, it
+   * is host memory; otherwise it is inline-eligible only if the
+   * cache entry says host. */
+  status = ucs_memtype_cache_lookup(buffer, length, &mem_info);
+  return (status == UCS_ERR_NO_ELEM) ||
+         ((status == UCS_OK) && (mem_info.type == UCS_MEMORY_TYPE_HOST));
\end{lstlisting}

The previous code was a blanket disable: once the memtype cache became
non-empty (which happens as soon as any CUDA buffer is registered), all
sends lost the inline-eligible fast path, including plain host buffers.
The fix restores the fast path for host memory while still excluding
GPU memory.

\subsection{PR \#11307 --- TCP as fallback wireup transport for RC/DC aliases}
\label{sec:pr_11307}

Fixes \href{https://github.com/openucx/ucx/issues/4794}{\#4794}. Head
commit \texttt{1554563a5} (test fix) + parent \texttt{a2ccbcb1f}
(production change). The production change is five lines in the
transport alias table:

\begin{lstlisting}[language=C]
-  { "rc_v",  { "rc_verbs", UCP_TL_AUX("ud_verbs"), NULL } },
-  { "rc_x",  { "rc_mlx5", UCP_TL_AUX("ud_mlx5"), NULL } },
+  { "rc_v",  { "rc_verbs", UCP_TL_AUX("ud_verbs"), UCP_TL_AUX("tcp"), NULL } },
+  { "rc_x",  { "rc_mlx5", UCP_TL_AUX("ud_mlx5"), UCP_TL_AUX("tcp"), NULL } },
   { "rc",    { "rc_mlx5", UCP_TL_AUX("ud_mlx5"), "rc_verbs",
-               UCP_TL_AUX("ud_verbs"), NULL } },
-  { "dc",    { "dc_mlx5", UCP_TL_AUX("ud_mlx5"), NULL } },
-  { "dc_x",  { "dc_mlx5", UCP_TL_AUX("ud_mlx5"), NULL } },
+               UCP_TL_AUX("ud_verbs"), UCP_TL_AUX("tcp"), NULL } },
+  { "dc",    { "dc_mlx5", UCP_TL_AUX("ud_mlx5"), UCP_TL_AUX("tcp"), NULL } },
+  { "dc_x",  { "dc_mlx5", UCP_TL_AUX("ud_mlx5"), UCP_TL_AUX("tcp"), NULL } },
\end{lstlisting}

The \texttt{UCP\_TL\_AUX(...)} macro marks a transport as
wireup-eligible only; it is not used for data traffic. UD remains the
preferred auxiliary (lower latency). TCP is only selected when UD is
unavailable, which happens precisely in the SoftRoCE / direct-connect
scenario that the r/HPC post described (see
Section~\ref{sec:reddit}).

The test-side fix (commit \texttt{1554563a5}) moves the TCP-skip
decision in \texttt{test\_ucp\_address\_v2.pack\_iface\_attrs} from a
macro-level condition into the test body, avoiding a segfault on
variants where \texttt{sender()} is not yet initialised when the skip
check fires.

\subsection{PR \#11308 --- Estimated bandwidth query on \texttt{ucp\_ep\_query}}
\label{sec:pr_11308}

Fixes \href{https://github.com/openucx/ucx/issues/6254}{\#6254}. Head
commit \texttt{ed3d69e91}. Adds \texttt{UCP\_EP\_ATTR\_FIELD\_ESTIMATED\_BW}
to the public \texttt{ucp\_ep\_attr\_t}, letting the caller specify a
local/remote memory type pair and receive the summed lane bandwidth.
For applications doing data-placement decisions (libscuffedrdma's own
WFA classifier is one), this removes the need to duplicate UCX's
internal lane-scoring code in user space. The implementation is a
read-only pass over \texttt{ucp\_ep\_config(ep)->key.rma\_bw\_lanes[]}
that is O(lanes) and cannot fail.

\subsection{PR \#11309 --- Symmetric traffic class propagation}
\label{sec:pr_11309}

Fixes \href{https://github.com/openucx/ucx/issues/10325}{\#10325}. Head
commit \texttt{8962d27d7}. Pure bit-field isolation fix in
\texttt{uct\_ib\_address\_unpack}: the traffic-class flag bit (bit 7)
was corrupting the RoCE version bits (5--6) because the unpack code
right-shifted without first masking. The fix adds an
\texttt{\textasciitilde UCT\_IB\_ADDRESS\_FLAG\_TRAFFIC\_CLASS} mask
before the shift, then uses the traffic class from the packed address
when the local interface has none configured. After this fix, both
endpoints end up with the same QoS class regardless of which one
configured it.

\section{External Validation: libmesh-rdma and the r/HPC Thread}
\label{sec:reddit}

On 2026-03-16, \texttt{autoscriptlabs/libmesh-rdma} appeared on GitHub:
an MIT-licensed C library (75~KB compiled) implementing RDMA networking
and a 55-function MPI subset for consumer GPU clusters \emph{without a
managed switch}. The author (user
\href{https://www.reddit.com/user/Ok-Pomegranate1314}{Ok-Pomegranate1314})
also posted to r/HPC describing exactly the problem PR~\#11307
addresses~\cite{reddit_directconnect_roce}:

\begin{quote}\small\itshape
Spent today fighting UCX's UD bootstrap on a direct-connect
ConnectX-7 ring (4x DGX Spark, no switch). You already know how this
goes: \texttt{ibv\_create\_ah()} needs ARP, ARP needs L2 resolution,
L2 resolution needs a subnet that both endpoints share or a switch
that routes between them. Without the switch, UCX dies in
\texttt{initial\_address\_exchange} and takes MPICH with it.
\end{quote}

The crucial measurement that validates the RC-over-TCP-bootstrap
architecture:

\begin{quote}\small\itshape
RC QPs don't need any of this. \texttt{ibv\_modify\_qp()} to RTR takes
the destination GID directly. No AH object. No ARP. No subnet
requirement beyond what the GID encodes. The firmware transitions the
QP just fine. 77~GB/s. 11.6~$\mu$s RTT. The transport layer works
perfectly on direct-connect RoCE. It's only the connection management
that's broken.
\end{quote}

The comment thread is equally illuminating because two independent
users tried to configure around the bug at the UCX layer:

\begin{itemize}
  \item \texttt{plan-bean} suggested \texttt{UCX\_TLS=\textasciicircum
    ud,ud:aux} to skip UD entirely. The OP rejected this because on
    RoCE, UCX's connection-manager path can still hit RDMACM, which in
    turn uses the same ARP-based address resolution. (This is true only
    when the application explicitly uses \texttt{UCP\_EP\_PARAM\_FIELD\_SOCK\_ADDR};
    the default \texttt{UCT}-auxiliary wireup path that PR~\#11307
    extends does not.)
  \item \texttt{MissionDependent4401} proposed a three-variable
    configuration: \texttt{UCX\_TLS=rc,tcp},
    \texttt{UCX\_RC\_VERBS\_TLS=rc}, \texttt{UCX\_TCP\_TLS=tcp}. This
    \emph{is} the environment-variable form of what PR~\#11307 makes
    the default behavior. It works, but requires the user to know that
    TCP can serve as a wireup auxiliary, and to set three variables
    correctly.
\end{itemize}

Read carefully, the thread is a real-time record of a contributor
abandoning UCX because the community consensus was that UCX could not
be made to work on direct-connect RoCE. PR~\#11307 is the minimal
in-tree fix that collapses those three environment variables into a
stock UCX build, and the 77~GB/s measurement from libmesh-rdma
demonstrates the underlying hardware path is already
production-quality once the bootstrap is sorted out. The PR is
therefore a contribution to a publicly-documented pain point with an
existence proof that it is worth fixing.

\paragraph{A footnote on RDMACM.} UCX supports two distinct wireup
paths: its native UCT-auxiliary wireup (which PR~\#11307 extends by
adding TCP as a valid auxiliary) and RDMACM-based wireup (a separate
code path used when the application supplies
\texttt{UCP\_EP\_PARAM\_FIELD\_SOCK\_ADDR} during
\texttt{ucp\_ep\_create}). PR~\#11307 targets the former. RDMACM-based
wireup remains unaffected; fixing that path would require either a
patch to \texttt{librdmacm} itself or an opt-out flag in UCX's
connection manager.

\section{Contrasting Production Designs}

\subsection{Apple MLX PR \#3094 (JACCL distributed backend)}

Apple's MLX framework added ring-topology distributed training support
via a custom ``JACCL'' backend~\cite{mlx_pr3094}. JACCL chose
\emph{CPU-mediated} collectives rather than async RDMA: sender and
receiver are tightly coupled via sequential send/recv operations to
prevent data corruption. Reported throughput is +124\% on 4-node
Devstral2 inference (135 $\to$ 303 tokens/sec on the prompt phase).

The tradeoff: the tight coupling eliminates entire classes of
race conditions (helpful for a research framework where correctness
trumps raw throughput), but it also serialises transfers and
prevents overlap. This is the opposite pole of the
async-CQ-with-priorities design UCX targets; libscuffedrdma's PMP
controller is specifically what makes async overlap manageable when
multiple priority classes share the same link.

\subsection{Pure Storage TurboQuant + FlashBlade}

Pure Storage's blog post~\cite{turboquant_flashblade} reports a
production deployment of TurboQuant KV cache compression paired with
FlashBlade object storage, achieving up to 10$\times$ faster KV cache
restores (5--6$\times$ compression, sustained 900--1000 MB/s per GPU on
DGX A100-40GB). The transport is RDMA plus GPUDirect Storage into GPU
memory; the compression layer is TurboQuant.

This is relevant to the thesis for two reasons. First, it demonstrates
that the TurboQuant-over-RDMA combination is already shipping in a
production storage product, not a research prototype. Second, it
clarifies where libscuffedrdma's contribution sits: FlashBlade moves
compressed bytes fast and NIXL abstracts the transport, but neither
classifies traffic by priority class or routes across a dual QP pool.
The WFA/PMP layer is the part that remains libscuffedrdma's.

\subsection{NVIDIA HPC-X}

NVIDIA HPC-X~\cite{nvidia_hpcx} is the reference downstream distribution
that bundles Open MPI, UCX, SHARP, UCC, and the NCCL UCX plug-in into a
single optimised toolkit. HPC-X ships with ConnectX firmware and
Quantum switch integrations; the SHARP component in particular is
NVIDIA-fabric-only because it offloads collective operations onto the
switch silicon. HPC-X is the correct choice for an all-NVIDIA
datacenter cluster.

\section{CI Status and Infrastructure Flakiness}

At time of writing, CI status across the six PRs is dominated by
infrastructure flakiness on the \texttt{swx-rain03} RoCE test node
(UMR CQ poll timeouts, UD endpoint timeouts, DC performance-envelope
timeouts). The same set of jobs fails across every PR regardless of
the code change, confirming that the failures are independent of the
patches under test. PR \#11304 is currently in a closed state after a
simultaneous force-push and close event on 2026-04-04; whether to
reopen it is deferred to Section~\ref{sec:next_steps} because the user
has not confirmed whether the close was intentional.

A separate survey of the ten most-recently-merged PRs in
\texttt{openucx/ucx} shows all of them merged with 151/152 checks
green, so the UCX maintainer workflow is to force-push-to-retrigger
until clean rather than to override red RoCE checks. The current
bottleneck is therefore scheduling (maintainer attention +
infrastructure stability) rather than code correctness.

No comments have been posted on any of the six PRs during this update
cycle: the author's preference is to keep the PR trackers focused on
code review and document context on the thesis side.

\section{Security Audit Findings}

Triggered by a memory-safety bug the author discovered in
libmesh-rdma's \texttt{mpi/mesh\_mpi\_pt2pt.c}
(\texttt{malloc(rhdr->length) + memcpy} using a network-supplied length
without validation against \texttt{wc.byte\_len}), this section
documents a parallel audit of both the six UCX PRs and the
libscuffedrdma middleware.

\subsection{UCX PRs --- all clean}

All six PRs were audited against the libmesh-rdma pattern and nine
related classes (integer overflow, signed/unsigned mismatch, NULL
dereference, buffer-over-read in address parsing, TCP bootstrap trust,
traffic class validation, pack/unpack symmetry, exception handling,
and \texttt{wc.byte\_len} validation). \textbf{No findings.} Confidence
is high across all six PRs. The changes are either configuration
threshold logic (\#11304), local state comparison (\#11305), send-side
application data only (\#11306), test fixes (\#11307), API additions
(\#11308), or bit-field isolation (\#11309). None introduce new
network-data-handling code paths.

One pre-existing concern, noted but \emph{not} in scope for these PRs:
\texttt{ucp\_am\_long\_first\_handler()} in
\texttt{src/ucp/core/ucp\_am.c} sums several network-supplied length
fields without an overflow check. This is unrelated to the thesis
PRs and would require a separate upstream patch.

\subsection{Middleware --- four high-severity findings}

Four places in \texttt{/home/nathan/ScuffedRDMA/middleware/} read a
network-supplied length and passed it to \texttt{recv()},
\texttt{frombuffer()}, or an array allocation without bounds checking.
All four are the direct Python analogue of the libmesh-rdma bug that
triggered the audit.

\begin{table}[h]
\centering\small
\begin{tabular}{@{}p{4.5cm}p{7cm}p{3cm}@{}}
\toprule
\textbf{Finding} & \textbf{Fix} & \textbf{Severity} \\
\midrule
\#1: \texttt{roce\_transport.py} JSON handshake unbounded recv &
Replaced entire JSON code path with fixed-size 64-byte binary
\texttt{QpInfo} via \texttt{rdma\_bootstrap} port &
Critical \\
\addlinespace
\#2: \texttt{vllm\_connector.py} unbounded \texttt{k\_len + v\_len} in
KV layer transfer &
Added \texttt{MAX\_KV\_LAYER\_BYTES = 500MB} constant and pre-recv
bounds check; wrapped \texttt{reshape()} in try/except on
\texttt{ValueError} &
Critical \\
\addlinespace
\#3: \texttt{sae\_steering.py} unbounded \texttt{nnz} in
\texttt{SparseVector.from\_bytes} &
Added \texttt{MAX\_SPARSE\_NNZ} / \texttt{MAX\_SPARSE\_DIM} bounds,
explicit buffer length check before \texttt{frombuffer}; narrowed the
\texttt{except} in \texttt{sync\_from\_remote} to only swallow
\texttt{ValueError} &
Critical \\
\addlinespace
\#4: \texttt{transport\_tcp\_sim.py} unbounded \texttt{resp\_len} in
async RDMA Read &
Added \texttt{MAX\_RDMA\_READ\_BYTES = 1GB} bound before
\texttt{readexactly} &
High \\
\bottomrule
\end{tabular}
\caption{Middleware findings and their remediation. All four are
  closed as of this update.}
\label{tab:sec_findings}
\end{table}

The attack model is trusted-peer on a management network, so these
findings are defense-in-depth rather than active-exploitation targets.
They matter for two reasons: a crashing handshake during a benchmark
run is a thesis-result-ruining event even without an adversary, and
the same code is on a path toward multi-tenant deployment
(Configurations D and E from Update 4) where crafted peers become a
real concern.

\section{Middleware Hardening: Adapting libmesh-rdma}

Alongside the audit, three new modules were added to
\texttt{middleware/}, each adapting a production-tested piece of
libmesh-rdma's C code into Python. Attribution is in a
\texttt{LICENSE-THIRD-PARTY} file at the middleware root. The MIT
license permits redistribution and adaptation with notice.

\subsection{\texttt{rdma\_gid\_discovery.py}}

Ports \texttt{mesh\_rdma\_find\_ipv4\_gid\_index()} and helpers from
\texttt{src/mesh\_rdma\_core.c}. Replaces the hardcoded
\texttt{gid\_index = 0} in \texttt{roce\_transport.py} with a scan of
the GID table that detects IPv4-mapped GIDs (10 zero bytes + 0xffff +
4 bytes of IPv4). When a \texttt{preferred\_ip} is given, the highest
matching index wins (RoCEv2 lives above RoCEv1 in the GID table).
Raises \texttt{RuntimeError} with a diagnostic summary if no
IPv4-mapped GID is found.

\subsection{\texttt{rdma\_bootstrap.py}}

Ports the \texttt{mesh\_rdma\_qp\_info\_t} struct and
\texttt{mesh\_rdma\_send\_handshake} / \texttt{mesh\_rdma\_accept\_handshake}
from \texttt{src/mesh\_rdma\_core.c}. The wire format is 64 bytes of
fixed-layout, network-byte-order struct: QPN, PSN, 16-byte GID, IPv4
address, GID index, MTU, and 28 bytes of padding for binary
compatibility with the C struct. The Python API is two functions
(\texttt{send\_handshake}, \texttt{accept\_handshake}) and one
dataclass (\texttt{QpInfo}) with \texttt{pack()} / \texttt{unpack()}
methods. A \texttt{\_recv\_exact} helper handles TCP's partial-read
semantics, which the previous JSON-based code ignored.

This module is what closes Security Finding \#1 from
Table~\ref{tab:sec_findings}: because the wire format is fixed size,
there is no length field for a crafted peer to manipulate, so the
unbounded-\texttt{recv} class of bug is gone by construction.

\subsection{\texttt{rdma\_qp\_state\_machine.py}}

Ports \texttt{mesh\_rdma\_create\_qp} and
\texttt{mesh\_rdma\_connect\_qp}. Wraps pyverbs's \texttt{QP} with a
\texttt{QueuePair} class that exposes
\texttt{to\_init() / to\_rtr() / to\_rts() / verify\_rts() / close()}.
Each transition is wrapped in a retry loop with exponential backoff
(10, 20, 40, 80, 160 ms), and \texttt{verify\_rts()} queries the QP
state after the final transition to catch silent failures.
\texttt{close()} forces a transition back to RESET before
\texttt{QP.close()} to avoid the \texttt{rdma\_rxe} zombie state
documented in Update 4 Section 9.5.

The key detail ported from libmesh-rdma is that \texttt{to\_rtr()} sets
the GRH fields (\texttt{dgid}, \texttt{sgid\_index}, \texttt{hop\_limit})
directly on the \texttt{AHAttr} rather than allocating an
\texttt{AH} object. This is what makes the wireup work on direct-connect
RoCE: the kernel never has to resolve an ARP entry because the
destination GID is embedded in the QP state machine transition.

\subsection{Integration with \texttt{roce\_transport.py}}

The existing \texttt{RoCETransport.connect()} method was retargeted to
use the three new modules: \texttt{find\_ipv4\_gid\_index} picks the
GID index, \texttt{send\_handshake} or \texttt{accept\_handshake}
exchanges \texttt{QpInfo} over TCP, and a \texttt{QueuePair} handles
the state machine. The old JSON handshake, the hardcoded
\texttt{gid\_index = 0}, and the inline RTR/RTS transitions are all
deleted. Net line change is approximately $-80$ LOC inline
$+30$ LOC call sites $+300$ LOC across the three new modules, plus a
test file.

\subsection{Testing}

\texttt{middleware/tests/test\_rdma\_bootstrap.py} contains 17 tests:
16 pure Python (wire format, handshake loopback, GID helpers, and
regression tests for the three non-handshake security findings) and
one hardware-dependent smoke test that verifies \texttt{QueuePair.to\_init()}
against real \texttt{rxe0}. The full RTR/RTS loopback test is guarded
so that it skips on SoftRoCE builds where the kernel refuses loopback
GRH transitions with EINVAL (a known rxe limitation that does not
apply to real ConnectX hardware). On the local testbed, 16 tests pass
and 1 skips.

\section{Next Steps and Open Questions}
\label{sec:next_steps}

Carried over from Update 4:

\begin{enumerate}
  \item Multi-threaded concurrent workload on SoftRoCE to demonstrate
        head-of-line blocking elimination.
  \item vLLM \texttt{KvConnector} integration for end-to-end
        prefill/decode over libscuffedrdma.
  \item Port the data path from pyverbs to C++ or Rust to remove the
        GIL bottleneck.
\end{enumerate}

New this update:

\begin{enumerate}
  \setcounter{enumi}{3}
  \item Confirm PR \#11307 behavior on a real direct-connect RoCE ring
        once ConnectX hardware is available, to verify the RDMACM
        footnote in Section~\ref{sec:reddit} holds.
  \item Decide whether PR \#11304 should be reopened; the close event
        was simultaneous with a force-push on 2026-04-04 and looks
        accidental, but this has not been confirmed.
  \item Decide whether to responsibly disclose the
        \texttt{mesh\_mpi\_pt2pt.c} vulnerability upstream to
        \texttt{autoscriptlabs/libmesh-rdma}. The project is MIT-licensed
        and in active development, so standard responsible-disclosure
        norms apply.
  \item Run the full \texttt{QueuePair} loopback RTR test on real
        ConnectX hardware to confirm the SoftRoCE skip is a rxe quirk
        and not a wrapper bug.
\end{enumerate}

\begin{thebibliography}{9}

\bibitem{libmesh_rdma}
autoscriptlabs, ``libmesh-rdma + libmesh-mpi: RDMA networking for
consumer GPU clusters --- no managed switch required.'' GitHub, 2026.
\url{https://github.com/autoscriptlabs/libmesh-rdma}

\bibitem{reddit_directconnect_roce}
Ok-Pomegranate1314, ``Custom MPI over RDMA for direct-connect RoCE ---
no managed switch, no UCX, no UD. 55 functions, 75 KB.'' r/HPC, 2026.
\url{https://www.reddit.com/r/HPC/comments/1ruzjw3/custom_mpi_over_rdma_for_directconnect_roce_no/}

\bibitem{mlx_pr3094}
ml-explore, ``MLX PR \#3094: JACCL distributed backend.'' GitHub, 2026.
\url{https://github.com/ml-explore/mlx/pull/3094}

\bibitem{turboquant_flashblade}
R.~Alvarez, ``Up to 10$\times$ faster KV cache restore: TurboQuant
meets FlashBlade.'' Pure Storage Blog, 2026.
\url{https://blog.purestorage.com/purely-technical/up-to-10x-faster-kv-cache-restore-turboquant-meets-flashblade/}

\bibitem{nvidia_hpcx}
NVIDIA, ``HPC-X: High-Performance Computing Software Toolkit.''
NVIDIA Networking, 2026.
\url{https://developer.nvidia.com/networking/hpc-x}

\bibitem{introl_gpu_virt}
Introl, ``GPU virtualization: maximizing utilization in multi-tenant
environments.'' Introl Blog, 2025.
\url{https://introl.com/blog/gpu-virtualization-maximizing-utilization-multi-tenant-environments}

\end{thebibliography}

\end{document}
